\documentclass[11pt,a4paper]{article}
\usepackage{fontspec}
\setmainfont{Latin Modern Roman}
\usepackage[spanish,es-nodecimaldot]{babel}
\usepackage[margin=2.2cm,headheight=14pt]{geometry}
\usepackage{graphicx}
\usepackage{xcolor}
\usepackage{tabularx}
\usepackage{array}
\usepackage{booktabs}
\usepackage{fancyhdr}
\usepackage{caption}
\usepackage[hidelinks]{hyperref}

\definecolor{azulJW}{HTML}{002D62}
\definecolor{grisJW}{HTML}{536D85}
\graphicspath{{images/}}
\setlength{\parindent}{0pt}
\setlength{\parskip}{6pt}
\renewcommand{\arraystretch}{1.2}
\captionsetup{font=small,labelfont=bf,justification=centering,hypcap=false}
\pagestyle{fancy}
\fancyhf{}
\fancyhead[L]{\small\color{grisJW}JaveWheels}
\fancyhead[R]{\small\color{grisJW}Informe de implementación}
\fancyfoot[C]{\small\thepage}
\renewcommand{\headrulewidth}{0.3pt}

\newcommand{\tituloPagina}[1]{%
  {\Large\bfseries\color{azulJW}#1}\par\medskip
}
\newcommand{\archivo}[1]{\texttt{\detokenize{#1}}}
\newcommand{\captura}[4]{%
  \begin{minipage}[t]{0.47\textwidth}
    \centering
    \includegraphics[width=\linewidth,height=#3\textheight,keepaspectratio]{#1}
    \captionof{figure}{#2}
    {\small\color{grisJW}#4\par}
  \end{minipage}%
}

\begin{document}

% Página 1: portada.
\thispagestyle{empty}
\vspace*{3cm}
{\large\color{grisJW}JaveWheels\par}
\vspace{0.7cm}
{\Huge\bfseries\color{azulJW}Implementación de navegación y vistas de JaveWheels\par}
\vspace{0.7cm}
{\large Aplicación Android de movilidad compartida\par}
\vspace{1.3cm}
\rule{\textwidth}{0.5pt}
\vspace{0.5cm}
{\large Informe de la implementación de las vistas por rol,\par
sus componentes y su validación en el emulador.\par}
\vfill
{\large Bogotá\par 1 de octubre de 2026\par}
\clearpage

% Página 2: alcance y estructura.
\tituloPagina{Alcance y organización de la implementación}

Se completó el recorrido visual del pasajero para consultar Wheels, abrir el
detalle de un recorrido, acceder a sus viajes y consultar una reserva confirmada.
La solución conserva la estructura de Kotlin, Jetpack Compose, Material 3 y
Navigation 3 del proyecto.

\textbf{Navegación implementada}

\begin{center}
\color{azulJW}
Inicio $\rightarrow$ Wheels disponibles $\rightarrow$ Detalle del Wheel\par
\smallskip
$\rightarrow$ Mis viajes / Reservas $\rightarrow$ Detalle de reserva
\end{center}

\emph{Buscar Wheels} abre los resultados. Cada botón \emph{Ver Wheel} transmite
el identificador del recorrido para mostrar sus datos. \emph{Solicitar cupo}
abre Mis viajes; la reserva confirmada permite consultar su detalle.
Los retrocesos retornan a la lista correspondiente. Una única barra inferior
selecciona Inicio en el flujo de Wheels y Mis viajes en el flujo de reservas.

\textbf{Archivos y recursos entregados}

{\small
\begin{tabularx}{\textwidth}{@{}>{\raggedright\arraybackslash}p{2.6cm}>{\raggedright\arraybackslash}X@{}}
\toprule
\textbf{Tipo} & \textbf{Archivos y responsabilidad} \\
\midrule
Pantallas nuevas &
\archivo{WheelsDisponiblesScreen.kt}, \archivo{DetalleWheelScreen.kt} y
\archivo{DetalleReservaScreen.kt}: resultados y detalles del pasajero. \\
Componente nuevo &
\archivo{ViajesPasajero.kt}: tarjetas y detalles. \archivo{AvatarUsuario.kt},
\archivo{SelectorRol.kt} y \archivo{MapaIlustrativo.kt}: componentes compartidos. \\
Recursos locales &
Logo, mapas y avatares se dibujan mediante componentes de Compose, sin PNG. \\
Navegación &
\archivo{Rutas.kt} y \archivo{JWNavHost.kt}: rutas, identificador del Wheel,
acciones de navegación y selección de la barra inferior. \\
Pantallas y botón &
\archivo{InicioScreen.kt}, \archivo{MisViajesScreen.kt} y
\archivo{Botones.kt}: acceso a resultados, pestañas internas y estilos
compartidos de botones. \\
Datos y modelos &
\archivo{Modelos.kt} y \archivo{DatosQuemados.kt}: modelos del pasajero y
listas de Wheels disponibles, reservas, historial y viajes publicados por el conductor. \\
Comentarios &
Cambios únicamente en comentarios de \archivo{InicioViewModel.kt},
\archivo{OnboardingViewModel.kt}, \archivo{RecuperarContrasenaViewModel.kt},
\archivo{Validaciones.kt} y \archivo{app/build.gradle.kts}. \\
Documentación &
\archivo{README.md}, ocho capturas en \archivo{docs/images} y la fuente
\archivo{Informe_Implementacion_Sebas.tex}. \\
\bottomrule
\end{tabularx}
}

\textbf{Reutilización y datos}

Las pantallas reutilizan el tema, los botones y el panel existentes.
\archivo{WheelPasajero} reúne conductor, ruta, vehículo, cupos, aporte y
recogida; \archivo{ReservaPasajero} agrega el estado y el detalle de recogida.
Los datos son locales: navegar no crea reservas ni modifica la disponibilidad.
\clearpage

% Página 3: capturas de Inicio y resultados.
\tituloPagina{Inicio y consulta de Wheels}

Inicio mantiene el cambio de rol y el panel del pasajero. Los resultados
muestran dos recorridos con conductor, ruta, horario, aporte, cupos y tipo de
recogida. La pantalla admite desplazamiento vertical y una fila horizontal
de filtros visuales. Cada tarjeta abre el detalle del Wheel seleccionado.

\vspace{0.4cm}
\noindent
\captura{inicio-pasajero.png}{Inicio del pasajero}{0.69}{Acceso a la búsqueda desde el panel principal.}
\hfill
\captura{wheels-disponibles.png}{Wheels disponibles}{0.69}{Tarjetas de recogida directa y Modo Buseta.}

\vfill
{\small\color{grisJW}Las capturas documentan la aplicación ejecutada en el emulador Android.}
\clearpage

% Página 4: capturas de los detalles.
\tituloPagina{Consulta del recorrido y de la reserva}

El detalle del Wheel presenta el conductor, la calificación, el vehículo,
los cupos y el aporte. El bloque de recogida incluye punto de encuentro,
distancia, tiempo a pie y hora aproximada. El detalle de reserva muestra el
estado confirmado y las acciones visuales de mensaje, compartir y cancelar.
Los mapas se dibujan con \archivo{MapaIlustrativo}, compartido con Inicio.
En los detalles se ajustan la altura y la etiqueta de recogida, sin botón de centrar.

\vspace{0.4cm}
\noindent
\captura{detalle-wheel.png}{Detalle del Wheel}{0.69}{Información del recorrido y acceso a Mis viajes.}
\hfill
\captura{detalle-reserva.png}{Detalle de reserva}{0.69}{Reserva confirmada y punto acordado de recogida.}

\vfill
{\small\color{grisJW}Las acciones pendientes muestran avisos breves y conservan los datos locales.}
\clearpage

% Página 5: pestañas y validación.
\tituloPagina{Mis viajes y resultado de la validación}
\small

Las pestañas \emph{Reservas} y \emph{Mis Wheels} separan las solicitudes activas
del historial en modo pasajero. La primera muestra una reserva confirmada y una solicitud
pendiente; la segunda, dos recorridos finalizados. El selector de rol se comparte con Inicio.

\noindent
\captura{mis-viajes-reservas.png}{Pestaña Reservas}{0.48}{Reserva confirmada y solicitud pendiente.}
\hfill
\captura{mis-viajes-wheels.png}{Pestaña Mis Wheels}{0.48}{Historial de recorridos finalizados.}

\medskip
\textbf{Compilación y comprobación del flujo}

La compilación de depuración del 1 de octubre de 2026 terminó con
\textbf{BUILD SUCCESSFUL}. Se ejecutó:

{\footnotesize
\begin{verbatim}
.\gradlew.bat :app:assembleDebug --max-workers=2 --no-configuration-cache
\end{verbatim}
}
El APK se genera en \archivo{app/build/outputs/apk/debug/app-debug.apk}.
En el emulador se comprobó el flujo del pasajero, los detalles, los retrocesos,
ambos roles y sus pestañas, el acceso al registro de conductor y el estado compartido
con Inicio y Perfil. Los detalles utilizan el mapa de Compose.

\textbf{Limitación de la validación}

\archivo{:app:lintDebug} no terminó: su análisis de comentarios entró en un
bucle en \archivo{BidirectionalTextDetector}. No se dispone de un resultado
de lint aprobado. El proyecto no contiene pruebas automatizadas existentes.

\textbf{Conclusión}

Las vistas por rol y la navegación quedan integradas y compiladas.
La búsqueda, los filtros, la gestión efectiva de cupos, la persistencia,
la mensajería y la ubicación en tiempo real requieren una implementación
posterior. Los estados actuales permiten revisar el comportamiento de las
pantallas con datos demostrativos.

\clearpage
\normalsize
\tituloPagina{Mis viajes del conductor y componentes compartidos}

El selector conserva el rol compartido con Inicio y Perfil. Si el usuario aún no
registró un vehículo, abre el flujo existente de registro. En modo conductor,
\emph{Programados} y \emph{Finalizados} muestran recorridos publicados por el usuario,
con horario, cupos ocupados y totales, aporte y estado. Al cambiar de rol se
selecciona la primera pestaña. Los botones de consulta muestran un aviso pendiente.

\medskip
\noindent
\captura{mis-viajes-conductor-programados.png}{Viajes publicados programados}{0.58}{Vista local del conductor.}
\hfill
\captura{mis-viajes-conductor-finalizados.png}{Viajes publicados finalizados}{0.58}{Historial local del conductor.}

\medskip
\archivo{LogoJaveWheels} dibuja un círculo amarillo y un carro azul mediante Compose.\par
\archivo{AvatarUsuario} dibuja el icono de persona
sobre un círculo celeste en Inicio, Perfil, tarjetas y detalles.\par
\archivo{MapaIlustrativo} conserva 300 dp en Inicio y utiliza 230 y 200 dp en
los detalles del Wheel y de la reserva, respectivamente.

\end{document}
